\documentclass[11pt]{article}
\usepackage[margin=1in]{geometry}
\usepackage{amsmath,amssymb,amsthm}
\usepackage[utf8]{inputenc}
\usepackage{listings}
\usepackage{hyperref}
\usepackage{xcolor}
\usepackage{textcomp}

\definecolor{leanblue}{rgb}{0.2,0.4,0.7}

\lstdefinelanguage{Lean}{
  keywords={theorem,lemma,def,axiom,structure,class,instance,where,by,sorry,simp,exact,intro,apply,have,let,fun,match,if,then,else,import,open,namespace,end,section,variable,inductive,noncomputable,abbrev,deriving},
  keywordstyle=\color{leanblue}\bfseries,
  comment=[l]{--},
  morecomment=[s]{/-}{-/},
  string=[b]",
  basicstyle=\ttfamily\small,
  breaklines=true,
  extendedchars=true,
  inputencoding=utf8,
  literate=
    {≥}{{$\geq$}}1
    {≤}{{$\leq$}}1
    {→}{{$\to$}}1
    {←}{{$\leftarrow$}}1
    {∧}{{$\wedge$}}1
    {∨}{{$\vee$}}1
    {¬}{{$\neg$}}1
    {∀}{{$\forall$}}1
    {∃}{{$\exists$}}1
    {⊆}{{$\subseteq$}}1
    {⊂}{{$\subset$}}1
    {∈}{{$\in$}}1
    {∉}{{$\notin$}}1
    {×}{{$\times$}}1
    {⊗}{{$\otimes$}}1
    {▸}{{$\blacktriangleright$}}1
    {⟨}{{$\langle$}}1
    {⟩}{{$\rangle$}}1
    {λ}{{$\lambda$}}1
    {α}{{$\alpha$}}1
    {β}{{$\beta$}}1
    {σ}{{$\sigma$}}1
    {θ}{{$\theta$}}1
    {ε}{{$\varepsilon$}}1
    {μ}{{$\mu$}}1
    {π}{{$\pi$}}1
    {Σ}{{$\Sigma$}}1
    {Π}{{$\Pi$}}1
    {▶}{{$\triangleright$}}1
    {⊔}{{$\sqcup$}}1
    {⊓}{{$\sqcap$}}1
    {⊥}{{$\bot$}}1
    {⊤}{{$\top$}}1
    {∅}{{$\emptyset$}}1
    {∪}{{$\cup$}}1
    {∩}{{$\cap$}}1
    {≠}{{$\neq$}}1
    {ℕ}{{$\mathbb{N}$}}1
    {₀}{{$_0$}}1
    {₁}{{$_1$}}1
    {₂}{{$_2$}}1
    {|>}{{$|{\!>}$}}2
    {·}{{$\cdot$}}1
    {—}{{---}}1
    {∘}{{$\circ$}}1
    {√}{{$\sqrt{}$}}1
    {∝}{{$\propto$}}1
    {≈}{{$\approx$}}1
    {═}{{=}}1
    {⟹}{{$\Longrightarrow$}}1,
}

\newtheorem{theorem}{Theorem}
\newtheorem{lemma}[theorem]{Lemma}
\newtheorem{definition}[theorem]{Definition}
\newtheorem{axiom_stmt}[theorem]{Axiom}

\title{Formal Verification: Coeffect Algebra}
\author{Zach Kelling, Lux Industries}
\date{December 2025}

\begin{document}
\maketitle

\begin{abstract}
We formalize the coeffect algebra that classifies what a computation needs from the world. Effects describe what a computation \emph{does}; coeffects describe what it \emph{needs}. The tensor product (list append) combines coeffects, with pure ($\emptyset$) as the identity.
\end{abstract}

\section{Introduction}
The coeffect algebra provides a formal framework for determining build reproducibility. A build is reproducible if and only if all its coeffects are reproducible (content-addressed I/O, TEE attestation). Non-reproducible coeffects (time, random, GPU, auth) contaminate the entire build.

\section{Definitions}
\begin{definition}[Coeffect]
Access types: pure, filesystem (CA/non-CA), network (CA/non-CA), time, random, GPU, auth, TEE.
\end{definition}

\begin{definition}[Coeffects]
A coeffect set (list of individual coeffects).
\end{definition}


\section{Theorems and Axioms}
\begin{theorem}[\texttt{empty\_is\_pure}]
Empty coeffects are pure.
\end{theorem}

\begin{theorem}[\texttt{pure\_is\_reproducible}]
Pure coeffects are reproducible.
\end{theorem}

\begin{theorem}[\texttt{tensor\_pure\_right}]
Right identity: $r \otimes \emptyset = r$.
\end{theorem}

\begin{theorem}[\texttt{tensor\_pure\_left}]
Left identity: $\emptyset \otimes r = r$.
\end{theorem}

\begin{theorem}[\texttt{tensor\_assoc}]
Associativity: $(a \otimes b) \otimes c = a \otimes (b \otimes c)$.
\end{theorem}

\begin{theorem}[\texttt{tensor\_pure\_pure}]
Pure tensor pure is pure: $\emptyset \otimes \emptyset = \emptyset$.
\end{theorem}

\begin{theorem}[\texttt{tensor\_reproducible}]
Reproducible tensor reproducible is reproducible.
\end{theorem}

\begin{theorem}[\texttt{time\_not\_reproducible}]
Time access is not reproducible.
\end{theorem}

\begin{theorem}[\texttt{random\_not\_reproducible}]
Random access is not reproducible.
\end{theorem}

\begin{theorem}[\texttt{ca\_filesystem\_reproducible}]
Content-addressed filesystem access is reproducible.
\end{theorem}

\begin{theorem}[\texttt{tee\_reproducible}]
TEE attestation is reproducible.
\end{theorem}


\section{Lean Source}
\lstinputlisting[language=Lean]{../../formal/lean/Build/Coeffect.lean}

\section{References}
\noindent Source code: \url{https://github.com/luxfi/proofs/blob/main/lean/Build/Coeffect.lean}\\
\noindent White papers: \url{https://papers.lux.network}

\end{document}
